\documentclass[11pt,a4paper]{article}
% ===================================================================
%  Vorlage: Lautschrift mit Unicode (xelatex + Charis SIL)
%  Kompilieren: xelatex vorlage_lautschrift_unicode.tex
%  Schrift: Charis SIL (kostenlos, software.sil.org/charis) – alternativ
%  Doulos SIL oder Gentium Plus. Lautzeichen direkt eintippen oder aus
%  ipa.typeit.org kopieren.
% ===================================================================
\usepackage{fontspec}
\usepackage[ngerman]{babel}
\usepackage[a4paper,margin=2cm]{geometry}
\usepackage{tikz}
\setmainfont{Charis SIL}
\newcommand\ipa[1]{[#1]}     % hier gibt es keinen Konflikt mit babel
\pagestyle{empty}
\setlength{\parindent}{0pt}

\begin{document}
{\Large\bfseries Das Vokalviereck}

\medskip
Mädchen \ipa{ˈmɛːtçən}, Abend \ipa{ˈaːbn̩t}, ich \ipa{ɪç}, Leute \ipa{ˈlɔʏ̯tə},
thing \ipa{θɪŋ}, un \ipa{œ̃}.

\begin{center}
\begin{tikzpicture}[x=1.6cm, y=1.6cm, every node/.style={font=\Large}]
  % Trapez: vorne oben, hinten oben, hinten unten, vorne unten
  \draw[thick] (0,3) -- (4,3) -- (4,0) -- (1.5,0) -- cycle;
  \draw[gray] (0.5,2) -- (4,2)  (1,1) -- (4,1)  (2,3) -- (2.75,0);
  % links vom Punkt ungerundet, rechts gerundet
  \foreach \x/\y/\v/\r in {0/3/i/y, 0.5/2/e/ø, 1/1/ɛ/œ, 4/3//u, 4/2//o, 4/1//ɔ}{
    \fill (\x,\y) circle (2pt);
    \node[left]  at (\x,\y) {\v};
    \node[right] at (\x,\y) {\r};}
  \node at (0.6,2.6) {ɪ}; \node at (1.0,2.6) {ʏ}; \node at (3.5,2.6) {ʊ};
  \node at (2.4,1.6) {ə}; \node at (2.3,0.45) {ɐ}; \node at (2.75,-0.3) {a};
  \node[font=\small, gray, above] at (0,3.15) {vorne};
  \node[font=\small, gray, above] at (4,3.15) {hinten};
\end{tikzpicture}
\end{center}
\end{document}
